\documentclass[a4paper,10pt]{book}
\usepackage[centertags]{amsmath}
\usepackage{amscd}
\usepackage{amsthm}
\usepackage{amssymb}
%\usepackage{array}
\usepackage{enumerate}
\usepackage{xypic}
%\usepackage[dvips]{graphics}
%\usepackage{moreverb}
\usepackage{multicol}

\usepackage[english,catalan]{babel}
%\usepackage{epsfig}
%\usepackage{rotating}
\usepackage[all]{xy}
\usepackage{color}
\usepackage{tikz}
%\usepackage{tabls} %per tenir m{\'e}s control sobre l'espai en les taules nom{\'e}s cal aquest paquet
\usepackage{indentfirst} %per comen\c{c}ar primer parragraf de seccio o capitol amb indent


\usepackage[T1]{fontenc}
\usepackage[sc]{mathpazo}
\linespread{1.1} % Palatino needs more leading (space between lines)

\usepackage[latin1]{inputenc}
\usepackage[twoside,bindingoffset=1cm]{geometry}


%%%%%%%%%%%%%%%%%%%%%%%%%%%%%%%%%%%%%%%%%%%%%%%%%%%%%%%%%%%%%%%%%%%%%%%%%
% theorem environments
%%%%%%%%%%%%%%%%%%%%%%%%%%%%%%%%%%%%%%%%%%%%%%%%%%%%%%%%%%%%%%%%%%%%%%%%%

\newtheorem{defn0}{Definition}[chapter]
\newtheorem{prop0}[defn0]{Proposition}
\newtheorem{thm0}[defn0]{Theorem}
\newtheorem{lemma0}[defn0]{Lemma}
\newtheorem{corollary0}[defn0]{Corollary}
\newtheorem{example0}[defn0]{Example}
\newtheorem{remark0}[defn0]{Remark}
\newtheorem{conjecture0}[defn0]{Conjecture}

\newenvironment{definition}{ \begin{defn0}}{\end{defn0}}
\newenvironment{proposition}{\bigskip \begin{prop0}}{\end{prop0}}
\newenvironment{theorem}{\bigskip \begin{thm0}}{\end{thm0}}
\newenvironment{lemma}{\bigskip \begin{lemma0}}{\end{lemma0}}
\newenvironment{corollary}{\bigskip \begin{corollary0}}{\end{corollary0}}
\newenvironment{example}{ \begin{example0}\rm}{\end{example0}}
\newenvironment{remark}{ \begin{remark0}\rm}{\end{remark0}}
\newenvironment{conjecture}{\begin{conjecture0}}{\end{conjecture0}}

\newcommand{\defref}[1]{Definition~\ref{#1}}
\newcommand{\propref}[1]{Proposition~\ref{#1}}
\newcommand{\thmref}[1]{Theorem~\ref{#1}}
\newcommand{\lemref}[1]{Lemma~\ref{#1}}
\newcommand{\corref}[1]{Corollary~\ref{#1}}
\newcommand{\exref}[1]{Example~\ref{#1}}
\newcommand{\secref}[1]{Section~\ref{#1}}
\newcommand{\remref}[1]{Remark~\ref{#1}}
\newcommand{\conjref}[1]{Conjecture~\ref{#1}}


%%%%%%%%%%%%%%%%%%%%%%%%%%%%%%%%%%%%%%%%%%%%%%%%%%%%%%%
% environments per la intro i pel resum en catal{\`a} %
%%%%%%%%%%%%%%%%%%%%%%%%%%%%%%%%%%%%%%%%%%%%%%%%%%%%%%%

%intro


\newtheorem{th1}{Theorem}
\newtheorem{prop1}[th1]{Proposition}
\newtheorem{cor1}[th1]{Corollary}
\newtheorem{lem1}[th1]{Lemma}
\newtheorem{de1}[th1]{Definition}
%catala

\newtheorem{th2}{Teorema}
\newtheorem{prop2}[th2]{Proposici{\'o}}
\newtheorem{cor2}[th2]{Corol{\cdot}lari}
\newtheorem{lem2}[th2]{Lema}

%%%%%%%%%%%%%%%%%%%%%%%%%%%%%%%%%%%%%%%%%%%%%%%%%%%%%%%%%%%%%%%%%%%%%%%%%%%
%%%% local definitions for this paper
%%%%%%%%%%%%%%%%%%%%%%%%%%%%%%%%%%%%%%%%%%%%%%%%%%%%%%%%%%%%%%%%%%%%%%%%%%%






%%%%%%%%%%%%%%%%%%%%%% aix{\`o} pels headings %%%%%%%%%%%%%%%%%%%%%%%%
\usepackage{fancyhdr}
\pagestyle{fancy}
\renewcommand{\chaptermark}[1]{\markboth{#1}{}}
\renewcommand{\sectionmark}[1]{\markright{\thesection\ #1}}
\fancyhf{} \fancyhead[LE,RO]{\bfseries\thepage}
\fancyhead[LO]{\bfseries\rightmark} \fancyhead[RE]{\bfseries\leftmark}

\def\paginaenblanc{\newpage%
\thispagestyle{empty}%
\vspace*{2cm}%
\newpage%
\thispagestyle{empty}%
}


%%%%%%%%%%%%%%%%%%%%%%%%%%%%%%%%%%%%%%%%%%%%%%%%%%%%%%%%%%%%%%%%%%%%%%%%%
% aux commands
%%%%%%%%%%%%%%%%%%%%%%%%%%%%%%%%%%%%%%%%%%%%%%%%%%%%%%%%%%%%%%%%%%%%%%%%%
%==========================================================================
% macros to support private authors' notes
%==========================================================================
\newif\ifprivate
\privatetrue
\def\xbar{\vskip0.09in\hrule\vskip0.06in}
\def\private#1{\ifprivate \xbar {\em #1} \xbar
\else \fi}
\def\huh{\ifprivate ??? \marginpar{\Huge ???}
\else \fi}
\def\???{\ifprivate {\bf {???}} \marginpar{\begin{center}{\Huge {\bf ?}}\end{center}}
\else \fi}
%\def\???{\ifprivate {\bf {???}} \marginpar{{\Huge {\bf ?}}}
%\else \fi}
\marginparsep1mm
\def\nota#1{\ifprivate  $\clubsuit$ \marginpar{\parbox[t]{2.4cm}{\begin{center}\tiny #1\end{center}}}
\else \fi}
\def\comment#1{\ifprivate \marginpar{\parbox[t]{2.4cm}{\begin{center}\tiny #1\end{center}}}
\else \fi}
%\def\nota#1{\ifprivate  $\clubsuit$ \marginpar{\parbox[t]{1.8cm}{\tiny #1}}
%\else \fi}
\def\privateeject{\ifprivate\eject\fi}
%\def\???{{\bf {???}} \marginpar{{\Huge {\bf ?}}} }
%%%%%%%%%%%%%%%%%%%%%%%%%%%%%%%%%%%%%%%%%%%%%%%%%%%%%%%%%%%%%%%%%%%%%%%%%%

%%%%%%%%%%%%%% definicions cap{\'\i}tol ""%%%%%%%%
%%%%%%%%%%%%%%%%%%%%%%%%%%%%%%%%%%
\def\av{\underline{a}}

\def\Zr{\mathbb Z^r} \def\Nr{\mathbb N^r}
\def\xv{\underline{x}}
\def\yv{\underline{y}}
\def\zv{\underline{z}}

\newcommand{\Proj}{\operatorname{Proj}}
\newcommand{\Spec}{\operatorname{Spec}}


\def\HiRIM#1{H^i_{\mathcal M}(#1)}

\newcommand{\im}{\operatorname{Im}}
\renewcommand{\ker}{\operatorname{Ker}}
\newcommand{\grade}{\operatorname{grade}}
\newcommand{\Ext}{\operatorname{Ext}}
\newcommand{\Hom}{\operatorname{Hom}}
\newcommand{\m}{\mathfrak m}
\newcommand{\n}{\mathfrak n}
\newcommand{\M}{\mathfrak m}
%\newcommand{\cL}{{\mathcal L}}
\newcommand{\cP}{{\mathcal P}}
\newcommand{\cE}{{\mathcal E}}
\newcommand {\ZZ}{\mathbb{Z}}
\newcommand{\cS}{{\mathcal S}}
\newcommand {\PP}{\mathbb{P}}

\def\p{\mathfrak{p}}
\def\q{\mathfrak{q}}
\DeclareMathOperator{\de}{deg}

\DeclareMathOperator{\Hl}{H}
 \DeclareMathOperator{\h}{h}

%%%%%%%%%%%%%%%%%%%%%%%%%%%%%%%%%%%%%%%%%%%%%%%%%%%%%%%%%%%%%%%%%%%%%%%%%
%%%%%%%%%%%%%%%%%%%%%%%%%%%%%%%%%%%%%%%%%%%%%%%%%%%%%%%%%%%%%%%%%%%%%%%%%
%%%%%%%%%%%%%%%%%%%%%%%%%%%%%%%%%%%%%%%%%%%%%%%%%%%%%%%%%%%%%%%%%%%%%%%%%
%%%%%%%%%%%%%%%%%%%%%%%%%%%%%%%%%%%%%%%%%%%%%%%%%%%%%%%%%%%%%%%%%%%%%%%%%
%%%%%%%%%%%%%%%%%%%%%%%%%%%%%%%%%%%%%%%%%%%%%%%%%%%%%%%%%%%%%%%%%%%%%%%%%
%%%%%%%%%%%%%%%%%%%%%%%%%%%%%%%%%%%%%%%%%%%%%%%%%%%%%%%%%%%%%%%%%%%%%%%%%
%%%%%%%%%%%%%%%%%%%%%%%%%%%%%%%%%%%%%%%%%%%%%%%%%%%%%%%%%%%%%%%%%%%%%%%%%
%%%%%%%%%%%%%%%%%%%%%%%%%%%%%%%%%%%%%%%%%%%%%%%%%%%%%%%%%%%%%%%%%%%%%%%%
%%%%%%%%%%%%%%%%%%%%%%%%%%%%%%%%%%%%%%%%%%%%%%%%%%%%%%%%%%%%%%%%%%%%%%%%
\begin{document}

\pagestyle{empty}

\begin{titlepage}
\begin{center}
\begin{figure}[htb]
\begin{center}
\includegraphics[width=6cm]{ub_color.pdf}
\end{center}
\end{figure}

\textbf{\LARGE Treball final de grau} \\
\vspace*{.5cm}
\textbf{\LARGE GRAU DE MATEM\`{A}TIQUES } \\
\vspace*{.5cm}
\textbf{\LARGE Facultat de Matem\`atiques i Inform\`atica\\ Universitat de Barcelona} \\
\vspace*{1.5cm}
\rule{16cm}{0.1mm}\\
\begin{Huge}
\textbf{AQUI EL T\'ITOL DEL TREBALL} \\
\end{Huge}
\rule{16cm}{0.1mm}\\

\vspace{1cm}

\begin{flushright}
\textbf{\LARGE Autor: el vostre nom}

\vspace*{2cm}

\renewcommand{\arraystretch}{1.5}
\begin{tabular}{ll}
\textbf{\Large Director:} & \textbf{\Large Dr. Nom tutor } \\
\textbf{\Large Realitzat a:} & \textbf{\Large  Departament....   } \\
 & \textbf{\Large (nom del departament)} \\
\\
\textbf{\Large Barcelona,} & \textbf{\Large \today }
\end{tabular}

\end{flushright}

\end{center}

\end{titlepage}


%%%%%%%%%%%%%%%%%%%%%%%%%%%%%%%%%%%%%%%%%%%%%%%%%%%%%%%%%%%%%%%%%%%%%%%%%
\newpage
\selectlanguage{english}
\pagenumbering{roman} \setcounter{page}{0}
\tableofcontents
\newpage \thispagestyle{empty}
%%%%%%%%%%%%%%%%%%%%%%%%%%%%%%%%%%%%%%%%%%%%%%%%%%%%%%%%%%%%%%%%%%%%%%%%%


%%%%%%%%%%%%%%%%%%%%% aix{\`o} pels headings %%%%%%%%%%%%%%%%%%%%%%%%
\pagestyle{fancy}
\addcontentsline{toc}{chapter}{Introduction}
\markboth{Introduction}{Introduction}
\newpage \thispagestyle{empty}


The main goal of this work is to prove the existence of.......



\vspace{3mm}

{\bf Notation:} We will work over an algebraically closed field $K$ of characteristic zero.

%%%%%%%%%%%%%%%%%%%%%%%%%%%%%%%%%%%%%%%%%%%%%%%%%%%%%%%%%%%%%%%%%%%%%%%%%
\mainmatter
\chapter[titol abreujat 1]{t\'{\i}tol del cap\'{\i}tol 1}
\section{Sub apartat del cap\'{\i}tol}

\begin{definition}  \label{ulrich}  \rm Given a projective scheme $X\subseteq \PP^n$ and a coherent sheaf $\cE$ on $X$, we say that
$\cE$ is an \emph{Ulrich sheaf} if  $\cE$ is an ACM sheaf and $\h^0(\cE_{init})=\deg(X)\mbox{rank}(\cE)$, i.e. the minimal
number of generators of Ulrich sheaves is as large as possible.
\end{definition}


Notice that in Definition \ref{ulrich} we do not require......

\vspace{3mm}


\begin{theorem}{\bf (Borel-Bott-Weil)}
\label{bottLine}
Assume that all entries of $\alpha+\rho$ are distinct.  Let $\sigma \in S_{n+1}$ be the unique permutation such
    that $\sigma(\alpha+\rho)$ is strictly decreasing. Then
\[ \Hl^m(F_c, L_{\alpha})= \left \{ \begin{array}{ll}
\Sigma^{\sigma(\alpha+\rho)-\rho} V^* & \mbox{if} \quad m=l(\sigma) \\
0 &\mbox{otherwise.}
 \end{array} \right. \]
If at least two entries of $\alpha + \rho $ coincide then $$
\Hl^{i}(F_c, L_{\alpha})= 0 \quad \mbox{for all integer } i.$$
\end{theorem}
\begin{proof}  See \cite{D}; Theorem 2.
\end{proof}




\begin{theorem}{\bf (Borel-Bott-Weil)}
\label{bott2} Let $ \beta=(\beta_1, \cdots, \beta_{p+1}) \in \ZZ^{p+1}$, $\alpha=(\alpha_1, \cdots, \alpha_{q-p}) \in \ZZ^{q-p}$ and $\gamma=(\gamma_1, \cdots, \gamma_{n-q}) \in \ZZ^{n-q}$
 be non-increasing integer sequences  and let $\mu=(\beta,\alpha, \gamma) \in \ZZ^{n+1}$ be their concatenation.

\begin{itemize}
\item[(a)] Assume that all entries of $\mu + \rho $ are distinct and let $\sigma \in S_{n+1}$ be the unique permutation such
    that $\sigma(\mu+\rho)$ is strictly decreasing. Then
\[ \Hl^m(F(p,q,n), \Sigma^{\beta} \cP \otimes \Sigma^{\alpha} \cS_2 \otimes \Sigma^{\gamma}\cS_1)= \left \{ \begin{array}{ll}
\Sigma^{\sigma(\mu+\rho)-\rho} V^* & \mbox{if} \quad m=l(\sigma) \\
0 &\mbox{otherwise.}
 \end{array} \right. \]
\item[(b)] If at least two entries of $\mu + \rho $ coincide then $$
\Hl^{i}(F(p,q,n), \Sigma^{\beta} \cP \otimes \Sigma^{\alpha} \cS_2 \otimes \Sigma^{\gamma}\cS_1)= 0 \quad \mbox{for all interger } i.$$
\end{itemize}
\end{theorem}
\begin{proof} It follows from Theorem \ref{bottLine}
\end{proof}

\chapter{t\'{\i}tol del cap\'{\i}tol 2}



%%%%%%%%%%%%%%%%%%%%%%%%%%%%%%%%%%%%%%%%%%%%%%%%%%%%%%%%%%%%%%%%%%%%%%%%%
\backmatter
\selectlanguage{english}
\addcontentsline{toc}{chapter}{Bibliography}
\bibliographystyle{amsplain}


\begin{thebibliography}{999}


   \bibitem{At} M.F. Atiyah, {\em Vector bundles over an elliptic curve}, Proc. London Math. Soc., {\bf 7}, no. 3, (1957),
    414--452.


 \bibitem{BGS}
    R. Buchweitz, G. Greuel  and F.O. Schreyer, {\em Cohen-Macaulay modules on hypersurface singularities, II},
     Invent. Math. , {\bf 88}, no.  1, (1987), 165--182.


\bibitem{D} M. Demazure, {\em A very simple proof of Bott's theorem}, Inventiones Math. {\bf 33} (1976) 271-272.




      \bibitem{EH} D. Eisenbud and J. Herzog,  {\em The classification of homogeneous {C}ohen-{M}acaulay rings of finite representation type}, Math. Ann. {\bf 280} (1988), 347--352.



\bibitem{ESW} D. Eisenbud, F.O. Schreyer  and J. Weyman,
    {\em Resultants and Chow forms via exterior syzygies},  J. of Amer. Math. Soc.,  {\bf 16}
    (2003), 537--579.

\bibitem{FM} D. Faenzi, F. Malaspina, {\em The CM representation type of homogeneous spaces}, Preprint 2014.


\end{thebibliography}



\end{document}
